% GENERATED FILE. Edit canonical lessons, then run npm run generate:books.
% canonical-source-hash: fnv1a64:89bb621092e299bd
% canonical-lessons: ZH-C484-yufang, ZH-C484-jiejiu, ZH-C484-yangcheng, ZH-C484-pilao, ZH-C484-dikangli

\chapter{Staying Well}
\label{ch:zh-staying-well}
\hlchaptermodality{484}

\begin{chapteropening}
\textbf{By the end of this chapter:} \emph{I can talk about preventing illness, giving up drinking, forming a habit, feeling worn out and fighting off colds.}

Complete the last lesson of chapter 484: I can talk about preventing illness, giving up drinking, forming a habit, feeling worn out and fighting off colds.
\end{chapteropening}

\section[yùfáng]{yùfáng — to prevent; to guard against}

\label{lesson:ZH-C484-yufang}



\begin{quote}
Before the new one: say the Mandarin for looking after one's health, then the Mandarin for chronic.
\end{quote}

\subsection*{You'll want to know: yùfáng}
\emph{yùfáng} — \textquotedblleft{}to prevent; to guard against\textquotedblright{}.

\begin{quote}
\emph{chángcháng xǐ shǒu kěyǐ yùfáng gǎnmào} — Washing your hands often helps prevent colds.
\end{quote}

Its characters come later; for now it is a word for the ear.

\begin{cousinweb}
\emph{yùfáng} is built from \emph{yù} 'beforehand' and \emph{fáng} 'to ward off', both inherited from Middle Chinese, \emph{yoH} and \emph{bjang} (Guangyun, in Baxter's transcription), which glosses the first 'in advance' and the second 'to guard against'.
\end{cousinweb}

\begin{grammarlens}[title={What you've built}]
One more word for saying how it was.
\end{grammarlens}

\subsection*{Your turn}
\begin{itemize}
\raggedright
  \item \emph{Say it:} \emph{yùfáng}
  \item \emph{Say it:} \emph{yùfáng}, once more
  \item \emph{Say it:} \emph{chángcháng xǐ shǒu kěyǐ yùfáng gǎnmào} — Washing your hands often helps prevent colds.
  \item \emph{Recall:} say the Mandarin for looking after one's health, then the Mandarin for chronic, then say \emph{yùfáng} again
\end{itemize}

\begin{tcolorbox}[breakable,colback=teal!4,colframe=teal!35!black,title={Before you move on}]
What does yùfáng mean? (\textquotedblleft{}To prevent; to guard against\textquotedblright{}.) Say it once more.
\end{tcolorbox}
\section[jièjiǔ]{jièjiǔ — to give up drinking}

\label{lesson:ZH-C484-jiejiu}



\begin{quote}
Before the new one: say the Mandarin for chronic, then the Mandarin for to guard against.
\end{quote}

\subsection*{You'll want to know: jièjiǔ}
\emph{jièjiǔ} — \textquotedblleft{}to give up drinking\textquotedblright{}.

\begin{quote}
\emph{yīshēng shuō tā yīnggāi jièjiǔ} — The doctor says he ought to give up drinking.
\end{quote}

Its characters come later; for now it is a word for the ear.

\begin{cousinweb}
\emph{jièjiǔ} is built from \emph{jiè} 'to guard against' and \zh{酒} 'alcohol', both inherited from Middle Chinese, \emph{keajH} and \emph{tsjuwX} (Guangyun, in Baxter's transcription), which glosses the first 'to be careful' and 'to be on guard' and the second as wine.
\end{cousinweb}

\begin{grammarlens}[title={What you've built}]
One more word for saying how it was.
\end{grammarlens}

\subsection*{Your turn}
\begin{itemize}
\raggedright
  \item \emph{Say it:} \emph{jièjiǔ}
  \item \emph{Say it:} \emph{jièjiǔ}, once more
  \item \emph{Say it:} \emph{yīshēng shuō tā yīnggāi jièjiǔ} — The doctor says he ought to give up drinking.
  \item \emph{Recall:} say the Mandarin for chronic, then the Mandarin for to guard against, then say \emph{jièjiǔ} again
\end{itemize}

\begin{tcolorbox}[breakable,colback=teal!4,colframe=teal!35!black,title={Before you move on}]
What does jièjiǔ mean? (\textquotedblleft{}To give up drinking\textquotedblright{}.) Say it once more.
\end{tcolorbox}
\section[yǎngchéng]{yǎngchéng — to form (a habit); to get into a habit}

\label{lesson:ZH-C484-yangcheng}



\begin{quote}
Before the new one: say the Mandarin for to guard against, then the Mandarin for to give up drinking.
\end{quote}

\subsection*{You'll want to know: yǎngchéng}
\emph{yǎngchéng} — \textquotedblleft{}to form (a habit); to get into a habit\textquotedblright{}.

\begin{quote}
\emph{wǒ yǎngchéng le zǎoshang pǎobù de xíguàn} — I've got into the habit of running in the morning.
\end{quote}

Its characters come later; for now it is a word for the ear.

\begin{cousinweb}
\emph{yǎngchéng} is built from \emph{yǎng} 'to raise' and \zh{成} 'to complete', both inherited from Middle Chinese, \emph{yangX} and \emph{dzyeng} (Guangyun, in Baxter's transcription), which glosses the first 'to rear' and the second 'to finish'.
\end{cousinweb}

\begin{grammarlens}[title={What you've built}]
One more word for saying how it was.
\end{grammarlens}

\subsection*{Your turn}
\begin{itemize}
\raggedright
  \item \emph{Say it:} \emph{yǎngchéng}
  \item \emph{Say it:} \emph{yǎngchéng}, once more
  \item \emph{Say it:} \emph{wǒ yǎngchéng le zǎoshang pǎobù de xíguàn} — I've got into the habit of running in the morning.
  \item \emph{Recall:} say the Mandarin for to guard against, then the Mandarin for to give up drinking, then say \emph{yǎngchéng} again
\end{itemize}

\begin{tcolorbox}[breakable,colback=teal!4,colframe=teal!35!black,title={Before you move on}]
What does yǎngchéng mean? (\textquotedblleft{}To form (a habit); to get into a habit\textquotedblright{}.) Say it once more.
\end{tcolorbox}
\section[píláo]{píláo — worn out; fatigued}

\label{lesson:ZH-C484-pilao}



\begin{quote}
Before the new one: say the Mandarin for to give up drinking, then the Mandarin for to get into a habit.
\end{quote}

\subsection*{You'll want to know: píláo}
\emph{píláo} — \textquotedblleft{}worn out; fatigued\textquotedblright{}.

\begin{quote}
\emph{wǒ gōngzuò le yì tiān gǎnjué hěn píláo} — I've worked all day, and I feel worn out.
\end{quote}

Its characters come later; for now it is a word for the ear.

\begin{cousinweb}
\emph{píláo} is built from \emph{pí} 'weary' and \emph{láo} 'toil', both inherited from Middle Chinese, \emph{bje} and \emph{law} (Guangyun, in Baxter's transcription), which glosses the first 'tired' and 'spent' and the second 'tired' and 'hard work'.
\end{cousinweb}

\begin{grammarlens}[title={What you've built}]
One more word for saying how it was.
\end{grammarlens}

\subsection*{Your turn}
\begin{itemize}
\raggedright
  \item \emph{Say it:} \emph{píláo}
  \item \emph{Say it:} \emph{píláo}, once more
  \item \emph{Say it:} \emph{wǒ gōngzuò le yì tiān gǎnjué hěn píláo} — I've worked all day, and I feel worn out.
  \item \emph{Recall:} say the Mandarin for to give up drinking, then the Mandarin for to get into a habit, then say \emph{píláo} again
\end{itemize}

\begin{tcolorbox}[breakable,colback=teal!4,colframe=teal!35!black,title={Before you move on}]
What does píláo mean? (\textquotedblleft{}Worn out; fatigued\textquotedblright{}.) Say it once more.
\end{tcolorbox}
\section[dǐkànglì]{dǐkànglì — resistance to illness}

\label{lesson:ZH-C484-dikangli}



\begin{quote}
Before the new one: say the Mandarin for to get into a habit, then the Mandarin for worn out.
\end{quote}

\subsection*{You'll want to know: dǐkànglì}
\emph{dǐkànglì} — \textquotedblleft{}resistance to illness\textquotedblright{}.

\begin{quote}
\emph{chángcháng duànliàn de rén dǐkànglì bǐ biérén hǎo} — People who exercise regularly fight off illness better than others.
\end{quote}

Its characters come later; for now it is a word for the ear.

\begin{cousinweb}
\emph{dǐkànglì} is built from \emph{dǐ} 'to push against', \emph{kàng} 'to resist' and \zh{力} 'strength', all inherited from Middle Chinese, \emph{tejX}, \emph{khangH} and \emph{lik} (Guangyun, in Baxter's transcription). The dictionary glosses the first 'to push' and lists the second under two readings, glossing \emph{khangH} as lifting with the hand.
\end{cousinweb}

\begin{grammarlens}[title={What you've built}]
One more word for saying how it was.
\end{grammarlens}

\subsection*{Your turn}
\begin{itemize}
\raggedright
  \item \emph{Say it:} \emph{dǐkànglì}
  \item \emph{Say it:} \emph{dǐkànglì}, once more
  \item \emph{Say it:} \emph{chángcháng duànliàn de rén dǐkànglì bǐ biérén hǎo} — People who exercise regularly fight off illness better than others.
  \item \emph{Recall:} say the Mandarin for to get into a habit, then the Mandarin for worn out, then say \emph{dǐkànglì} again
\end{itemize}

\begin{tcolorbox}[breakable,colback=teal!4,colframe=teal!35!black,title={Before you move on}]
What does dǐkànglì mean? (\textquotedblleft{}Resistance to illness\textquotedblright{}.) Say it once more.
\end{tcolorbox}
